\chapter{Cơ sở lý thuyết}
\section{Xử lý ngôn ngữ tự nhiên}
Xử lý ngôn ngữ tự nhiên, hay NLP, là lĩnh vực nghiên cứu cách máy tính xử lý, hiểu và sinh ngôn ngữ của con người \cite{jurafsky2023speech}. Trong đề tài này, NLP xuất hiện ở nhiều tầng: làm sạch review, biểu diễn văn bản bằng \TFIDF{}, phân loại sentiment, biểu diễn câu bằng embedding, truy xuất review liên quan và tạo câu trả lời dựa trên bằng chứng.

Văn bản tiếng Việt có một số đặc thù. Từ có thể gồm nhiều âm tiết cách nhau bằng khoảng trắng, ví dụ “giao hàng”, “chất lượng”, “không trả lời”. Người dùng thường viết tắt, thiếu dấu, dùng emoji hoặc lặp ký tự. Review TMĐT còn có nhiều câu rất ngắn như “ok”, “tạm được”, “giao lâu”, hoặc câu dài nhưng nhiễu do người dùng viết để nhận xu. Vì vậy tiền xử lý và lựa chọn mô hình cần chú ý đến nhiễu dữ liệu.

\section{Tiền xử lý văn bản}
Tiền xử lý văn bản là bước chuyển dữ liệu thô thành dạng có thể dùng cho mô hình. Với review TMĐT, các bước phổ biến gồm chuẩn hóa Unicode, loại bỏ HTML, xóa khoảng trắng thừa, kiểm tra giá trị rỗng, chuẩn hóa rating và xử lý duplicate. Trong đề tài, Module 1 thực hiện các bước này trước khi tạo tập train/valid/test cho sentiment và corpus cho RAG.

Tiền xử lý không nên làm mất thông tin quan trọng. Ví dụ, từ “không” là tín hiệu phủ định quan trọng; nếu xóa tùy tiện, câu “không tốt” có thể bị hiểu thành “tốt”. Các từ như “lỗi”, “chậm”, “thiếu”, “sai”, “hỏng”, “bể”, “vỡ”, “kém”, “tệ” cũng là tín hiệu complaint quan trọng cho retrieval.

\section{Phân tích cảm xúc}
Phân tích cảm xúc là bài toán xác định thái độ hoặc cảm xúc được thể hiện trong văn bản. Trong đề tài này, sentiment được chia thành ba lớp: negative, neutral và positive. Review “Shop giao sai màu, nhắn tin không trả lời” thường là negative. Review “Sản phẩm dùng ổn, giao nhanh” thường là positive. Review “Tạm được, chưa dùng lâu nên chưa biết” có thể là neutral.

Bài toán sentiment classification trong đề tài nhận input là một review và output là nhãn sentiment. Đây là bài toán phân loại văn bản có giám sát, vì dữ liệu train có nhãn. Nhãn được ánh xạ từ rating: 1--2 sao là negative, 3 sao là neutral, 4--5 sao là positive. Cách ánh xạ này đơn giản và phù hợp với dữ liệu có rating, nhưng có thể tạo nhiễu vì rating không luôn phản ánh đúng nội dung review.

\section{\TFIDF{}}
\TFIDF{} là phương pháp biểu diễn văn bản bằng vector dựa trên tần suất từ \cite{salton1988term}. TF đo một từ xuất hiện nhiều hay ít trong một văn bản. IDF giảm trọng số của các từ xuất hiện ở quá nhiều văn bản. Ví dụ, trong review TMĐT, từ “sản phẩm” có thể xuất hiện rất nhiều nên ít phân biệt. Ngược lại, các từ “giao sai”, “móp”, “không trả lời”, “hỏng” có thể mang nhiều thông tin hơn.

Với \TFIDF{}, câu “shop giao sai màu” được chuyển thành vector sparse, trong đó các chiều tương ứng với từ hoặc n-gram như “shop”, “giao”, “sai”, “màu”, “giao sai”. Hệ thống dùng n-gram 1--2 trong mô hình tốt nhất, tức là xét cả từ đơn và cụm hai từ. Điều này quan trọng vì cụm “không tốt” khác nhiều so với từ “tốt” đứng riêng.

\section{Logistic Regression và Linear SVM}
Logistic Regression là mô hình tuyến tính dùng để ước lượng xác suất một mẫu thuộc về từng lớp. Với text classification, input là vector \TFIDF{}, output là xác suất hoặc điểm cho các lớp negative, neutral, positive. Mô hình học trọng số cho từng từ/cụm từ. Nếu từ “tệ”, “lỗi”, “không trả lời” thường xuất hiện trong review negative, trọng số của chúng cho lớp negative sẽ lớn.

Linear SVM là mô hình phân loại tuyến tính tìm siêu phẳng phân tách các lớp với biên lớn \cite{cortes1995support}. Trong dữ liệu text, vector \TFIDF{} thường có số chiều lớn và sparse. Linear SVM phù hợp với dạng dữ liệu này vì nó hoạt động tốt trên không gian chiều cao và thường cho kết quả mạnh trong text classification \cite{fan2008liblinear}. Trong thực nghiệm của đề tài, Linear SVM đạt kết quả cao nhất trong các baseline đã thử: test accuracy khoảng 0,9329 và test \MacroFOne{} khoảng 0,9098.

\section{Majority baseline và các metric}
Majority baseline là mô hình rất đơn giản: luôn dự đoán lớp xuất hiện nhiều nhất trong tập train. Mô hình này không có giá trị ứng dụng cao, nhưng là mốc tối thiểu để so sánh. Nếu mô hình phức tạp không vượt được majority baseline, mô hình đó không hữu ích.

Accuracy là tỷ lệ dự đoán đúng trên toàn bộ mẫu. Precision của một lớp đo trong các mẫu được dự đoán là lớp đó, bao nhiêu mẫu thật sự đúng. Recall đo trong các mẫu thật sự thuộc lớp đó, mô hình tìm đúng được bao nhiêu. F1-score là trung bình điều hòa của precision và recall. Macro-F1 tính F1 trung bình đều cho từng lớp, không phụ thuộc lớp nhiều hay ít mẫu. Weighted-F1 cũng trung bình theo lớp nhưng có trọng số theo số mẫu. Trong dữ liệu lệch lớp, \MacroFOne{} quan trọng vì nó không để lớp lớn che khuất lớp nhỏ.


\section{Các công thức nền tảng}
Các metric và thuật toán trong báo cáo có thể tóm tắt bằng một số công thức chính. \TFIDF{} được viết ở dạng khái niệm:
\begin{equation}
\mathrm{tfidf}(t,d)=\mathrm{tf}(t,d)\times \log\frac{N}{\mathrm{df}(t)}.
\end{equation}
Công thức này làm nổi bật các từ/cụm từ phân biệt như "giao sai", "hàng lỗi", "không trả lời", đồng thời giảm trọng số từ quá phổ biến như "sản phẩm" \cite{salton1988term,manning2008ir}.

Với một lớp sentiment, precision, recall và F1 được tính như sau:
\begin{align}
\mathrm{Precision} &= \frac{TP}{TP+FP},\\
\mathrm{Recall} &= \frac{TP}{TP+FN},\\
F_1 &= \frac{2\times \mathrm{Precision}\times \mathrm{Recall}}{\mathrm{Precision}+\mathrm{Recall}}.
\end{align}
Macro-F1 tính trung bình đều theo lớp:
\begin{equation}
\mathrm{MacroF1}=\frac{1}{C}\sum_{c=1}^{C}F_{1,c}.
\end{equation}
Vì neutral chỉ chiếm 871/5.162 review sau dedup, \MacroFOne{} phù hợp hơn accuracy khi chọn mô hình.

Với embedding đã normalize, inner product tương đương cosine similarity:
\begin{equation}
\cos(q,d)=\frac{q\cdot d}{\|q\|\|d\|}=q\cdot d \quad \text{khi } \|q\|=\|d\|=1.
\end{equation}
Đây là lý do Module 4 dùng FAISS \code{IndexFlatIP} cho embedding \MultilingualE{} \cite{johnson2019billion,wang2022e5}.

Linear SVM tối ưu biên phân tách lớn. Ở mức khái niệm, binary SVM có thể viết:
\begin{equation}
\min_{w,b} \frac{1}{2}\|w\|^2 + C\sum_i \max(0,1-y_i(w^Tx_i+b)).
\end{equation}
Dạng triển khai multi-class trong thư viện giúp mô hình phù hợp với vector \TFIDF{} sparse \cite{cortes1995support,fan2008liblinear,pedregosa2011scikit}.

BM25 và RRF được Module 6 dùng để nâng cấp retrieval. BM25 có dạng:
\begin{equation}
\mathrm{BM25}(q,d)=\sum_{t\in q}\mathrm{IDF}(t)\frac{f(t,d)(k_1+1)}{f(t,d)+k_1(1-b+b\frac{|d|}{\mathrm{avgdl}})}.
\end{equation}
RRF gộp nhiều ranking bằng:
\begin{equation}
\mathrm{RRF}(d)=\sum_i \frac{1}{k+\mathrm{rank}_i(d)}.
\end{equation}
Các công thức này giải thích vì sao BM25 bắt keyword tốt, dense retrieval bắt ngữ nghĩa tốt, còn hybrid có thể tận dụng cả hai \cite{robertson2009probabilistic,cormack2009rrf}.

RAG có thể tóm tắt thành Retrieve $\rightarrow$ Generate \cite{lewis2020rag}. Trong hệ thống này, Generate là template-based, không phải LLM tự do. Các hướng nâng cấp như transformer, PhoBERT, XLM-R, Sentence-BERT, ColBERT hoặc reranker neural được dùng như bối cảnh lý thuyết, chưa thay đổi kết luận thực nghiệm hiện tại \cite{devlin2019bert,vaswani2017attention,conneau2020xlmr,nguyen2020phobert,reimers2019sentencebert,nogueira2019passage,khattab2020colbert}. Các khảo sát sentiment analysis cũng cho thấy label noise và dữ liệu theo domain là vấn đề trung tâm của bài toán phân tích cảm xúc \cite{pang2008opinion,liu2012sentiment}.

\section{Embedding, dense retrieval và FAISS}
Embedding là vector dense biểu diễn ý nghĩa của văn bản. Khác với \TFIDF{} sparse vector dựa trên từ khóa, embedding cố gắng đặt các câu có nghĩa gần nhau ở vị trí gần nhau trong không gian vector. Ví dụ, “giao hàng lâu” và “ship trễ” có thể có embedding gần nhau dù không trùng từ hoàn toàn.

Dense retrieval là quá trình tìm tài liệu liên quan bằng embedding. Mỗi review được mã hóa thành vector. Query của người dùng cũng được mã hóa thành vector. Hệ thống tìm các review có vector gần query nhất. Hướng tiếp cận dense retrieval đã được dùng rộng rãi trong question answering và retrieval hiện đại \cite{karpukhin2020dense}. Hệ thống dùng FAISS \code{IndexFlatIP} \cite{johnson2019billion}. Vì embedding được normalize, inner product gần tương đương cosine similarity.

Hệ thống dùng \code{intfloat/multilingual-e5-small}, thuộc họ E5 text embedding \cite{wang2022e5}. Theo quy ước E5, document được mã hóa với prefix \code{passage: } và query được mã hóa với prefix \code{query: }. Prefix giúp mô hình hiểu vai trò của chuỗi đầu vào trong bài toán retrieval.

\section{RAG và \EvidenceGrounded{} generation}
RAG là mô hình kiến trúc gồm hai bước: retrieve và generate \cite{lewis2020rag}. Bước retrieve tìm tài liệu liên quan. Bước generate tạo câu trả lời dựa trên tài liệu đó. Trong các hệ thống hiện đại, bước generate thường dùng LLM. Tuy nhiên, đề tài này dùng generator dạng template-based để đảm bảo ổn định, không cần API key và không bịa ngoài dữ liệu.

Điểm quan trọng là câu trả lời phải \EvidenceGrounded{}. Nghĩa là mỗi ý chính cần dựa trên review được truy xuất. Nếu evidence yếu, hệ thống phải trả lời thận trọng, ví dụ “Chưa đủ bằng chứng rõ ràng từ các review được truy xuất”.

\section{BM25, RRF, reranking và ABSA}
BM25 là phương pháp truy xuất dựa trên keyword \cite{robertson2009probabilistic}. Nó mạnh khi query và document có từ giống nhau, ví dụ query “sai size” và review có cụm “giao sai size”. Dense retrieval lại mạnh khi query và review khác từ nhưng gần nghĩa. Vì vậy Module 6 kết hợp BM25 và dense retrieval.

RRF, hay Reciprocal Rank Fusion, là phương pháp gộp nhiều danh sách ranking \cite{cormack2009rrf}. Một document xuất hiện ở thứ hạng cao trong cả BM25 và dense sẽ được đẩy lên. Reranking là bước sắp xếp lại candidates sau retrieval. Module 6 có optional CrossEncoder reranker; nếu không load được, hệ thống fallback về lexical reranker.

Aspect-based sentiment analysis phân tích cảm xúc theo từng khía cạnh, ví dụ chất lượng, giao hàng, đóng gói, giá, dịch vụ shop, size/màu \cite{pontiki2014semeval}. Module 6 mới là MVP rule-based, chưa phải mô hình ABSA có giám sát. Nó dùng từ khóa để phát hiện aspect và gán sentiment aspect theo lexicon đơn giản. Cách này minh bạch trong phạm vi thực nghiệm, nhưng chưa xử lý tốt mỉa mai, phủ định phức tạp hoặc lỗi chính tả.

\section{Phân biệt các thành phần}
Sentiment classification trả lời câu hỏi: review này là tiêu cực, trung lập hay tích cực. Review analytics trả lời câu hỏi: toàn bộ dữ liệu có phân phối sentiment/rating như thế nào. Retrieval trả lời câu hỏi: review nào liên quan đến query. RAG trả lời câu hỏi: dựa trên các review liên quan, có thể tóm tắt insight gì và bằng chứng là gì. Phân biệt rõ các thành phần này giúp tránh hiểu nhầm rằng RAG thay thế sentiment model, hoặc sentiment model đang phân tích câu hỏi RAG như một review.

\section{Đặc thù tiếng Việt trong review thương mại điện tử}
Tiếng Việt trong review thương mại điện tử không giống văn bản chuẩn trong sách báo. Người dùng có thể viết không dấu, viết tắt, dùng tiếng lóng, emoji hoặc gõ sai chính tả. Ví dụ, “giao hang lau qua”, “shop rep cham”, “hang loi ko dung dc” đều là những biểu thức không chuẩn nhưng rất quan trọng với bài toán complaint mining. Nếu hệ thống chỉ dựa vào từ chuẩn, các tín hiệu này có thể bị bỏ sót.

Một thách thức khác là khoảng trắng trong tiếng Việt không luôn tương ứng với ranh giới từ. Cụm “giao hàng” gồm hai âm tiết nhưng biểu diễn một khái niệm. Khi dùng \TFIDF{} ở mức token split đơn giản, “giao” và “hàng” có thể bị tách. Vì vậy n-gram 1--2 hữu ích vì giữ được cụm “giao hàng”, “chất lượng”, “không tốt”, “sai màu”. Với retrieval, query expansion cũng giúp nối các biến thể như “ship”, “vận chuyển”, “giao lâu”, “giao chậm”.

Review TMĐT còn có hiện tượng rating và nội dung không khớp. Người dùng có thể cho 5 sao nhưng viết “giao hơi lâu”, hoặc cho 1 sao vì vấn đề shipper chứ không phải sản phẩm. Đây là lý do báo cáo không xem nhãn rating-derived sentiment là tuyệt đối đúng. Trong phân tích lỗi, các trường hợp lệch rating-text cần được nêu rõ để thể hiện tính trung thực khoa học.

\section{Ví dụ trực quan về \TFIDF{} và n-gram}
Giả sử có ba review: “giao hàng nhanh”, “giao hàng chậm”, “sản phẩm tốt”. Nếu chỉ nhìn từ đơn, hai review đầu đều có “giao” và “hàng”, nhưng khác nhau ở “nhanh” và “chậm”. Nếu dùng bigram, hệ thống có thêm đặc trưng “giao hàng”, “hàng nhanh”, “hàng chậm”. Với sentiment, bigram “không tốt” đặc biệt quan trọng vì từ “tốt” đứng riêng mang nghĩa tích cực, nhưng “không tốt” lại là tiêu cực.

Trong đề tài, \TFIDF{} với n-gram 1--2 và Linear SVM đạt kết quả tốt. Điều này không có nghĩa \TFIDF{} hiểu ngữ nghĩa sâu như transformer, nhưng cho thấy trong dữ liệu review ngắn, nhiều tín hiệu sentiment nằm ở từ khóa và cụm từ ngắn. Vì vậy baseline cổ điển là lựa chọn hợp lý cho phạm vi thực nghiệm: đơn giản, dễ tái lập và có thể diễn giải.

\section{Vì sao chọn validation \MacroFOne{}}
Nếu chỉ chọn model theo accuracy, mô hình có thể ưu tiên lớp lớn. Trong dữ liệu sentiment của đề tài, positive và negative nhiều hơn neutral. Một mô hình dự đoán tốt positive/negative nhưng bỏ qua neutral vẫn có accuracy khá cao. Macro-F1 tính đều cho ba lớp, nên model phải quan tâm cả neutral. Đây là lý do hệ thống chọn model bằng validation \MacroFOne{} thay vì chỉ accuracy.

Macro-F1 cũng phù hợp với mục tiêu học thuật vì nó kiểm tra mức độ cân bằng giữa các lớp. Việc chỉ báo cáo accuracy có thể che khuất lỗi ở lớp ít mẫu. Per-class F1 trong Chương 6 cho thấy neutral là lớp yếu hơn, từ đó tạo cơ sở cho phần phân tích lỗi và hướng phát triển.

\section{Niềm tin, bằng chứng và hallucination trong RAG}
Một hệ thống RAG tốt không chỉ cần trả lời đúng, mà còn cần cho biết câu trả lời dựa trên đâu. Trong đề tài này, citation đóng vai trò ràng buộc. Mỗi ý trong answer phải khớp với evidence card. Nếu evidence yếu, hệ thống phải nói thận trọng. Đây là cách giảm hallucination trong điều kiện không dùng LLM thật.

Nếu dùng LLM tự do mà không kiểm soát, model có thể viết câu trả lời nghe tự nhiên nhưng không có review nào chứng minh. Với bài toán phân tích review TMĐT, điều đó có thể dẫn đến quyết định sai. Do đó hệ thống chọn template-based generation: kém linh hoạt hơn nhưng minh bạch hơn. Đây là trade-off có chủ đích.

\section{Đánh giá retrieval khác gì đánh giá classifier}
Classifier có nhãn test rõ ràng: mỗi review có sentiment label. Retrieval khó hơn vì relevance phụ thuộc query. Một review có thể liên quan đến query “đóng gói” nhưng không liên quan đến query “size”. Vì vậy retrieval cần bộ query đánh giá và nhãn relevance theo từng cặp query-review. Manual Precision@k trong đề tài là một bước đầu để đo relevance này.

Điểm score của FAISS chỉ cho biết vector gần nhau trong không gian embedding. Nó không đảm bảo review trả lời đúng câu hỏi. Ví dụ query hỏi “phàn nàn về giao hàng chậm” có thể lấy review “giao hàng nhanh” vì cùng chủ đề giao hàng. Do đó sentiment filter, query expansion, keyword rerank và Module 6 hybrid retrieval đều nhằm giảm lỗi cùng chủ đề nhưng sai intent.

